% Generated draft: requires mathematical and source review.
\documentclass[11pt]{article}
\usepackage[margin=25mm]{geometry}
\usepackage{fontspec}
\setmainfont{Latin Modern Roman}
\usepackage{amsmath,amssymb,mathtools,unicode-math}
\setmathfont{Latin Modern Math}
\usepackage{enumitem,needspace,xurl,hyperref}
\setlength{\parindent}{0pt}
\setlength{\parskip}{5pt}
\setlength{\emergencystretch}{3em}
\newcommand{\N}{\mathbb{N}}
\newcommand{\Z}{\mathbb{Z}}
\newcommand{\Q}{\mathbb{Q}}
\newcommand{\R}{\mathbb{R}}
\newcommand{\C}{\mathbb{C}}
\newcommand{\F}{\mathbb{F}}
\newcommand{\A}{\mathbb{A}}
\newcommand{\PP}{\mathbb{P}}
\newcommand{\E}{\mathbb{E}}
\newcommand{\eps}{\varepsilon}
\newcommand{\dd}{\,\mathrm{d}}
\begin{document}
\section*{UnsolvedMath ID: 801197}
\textbf{Source problem code:} catalogue entry
\subsection{Definitions and mathematical statement}
Use the usual convention that a graph is undirected and simple; the supplied record does not explicitly specify this convention. Let \(G=(V,E)\) be a finite graph and put \(n=|V|\). The graph \(G\) is \emph{claw-free} if it has no induced subgraph isomorphic to \(K_{1,3}\), the graph consisting of a central vertex adjacent to three pairwise nonadjacent vertices.

A proper coloring is a map \(\kappa:V\to\{1,2,3,\ldots\}\) such that \(\kappa(u)\ne\kappa(v)\) whenever \(\{u,v\}\in E\). For commuting indeterminates \(x=(x_1,x_2,\ldots)\), define the chromatic symmetric function by
\[
X_G(x)=\sum_{\kappa\text{ proper}}\prod_{v\in V}x_{\kappa(v)}.
\]
This is a homogeneous symmetric function of degree \(n\). For a partition \(\lambda\vdash n\), let \(s_\lambda(x)\) be the Schur function, defined as the sum of the monomials \(\prod_{b\in\lambda}x_{T(b)}\) over all fillings \(T\) of the Young diagram of \(\lambda\) by positive integers that are weakly increasing along rows and strictly increasing down columns. Write the unique Schur expansion as
\[
X_G(x)=\sum_{\lambda\vdash n}a_\lambda(G)s_\lambda(x).
\]
The function \(X_G\) is \emph{Schur-positive}, or \(s\)-positive, if \(a_\lambda(G)\in\Z_{\ge0}\) for every \(\lambda\vdash n\). For \(n=0\), use the empty partition and \(X_G=s_{\varnothing}=1\).

\textbf{Conjecture (attributed to Vesselin Gasharov in the supplied record).} For every finite claw-free graph \(G\), its chromatic symmetric function is Schur-positive; equivalently,
\[
\forall\,G\text{ finite and claw-free},\quad
\forall\,\lambda\vdash |V(G)|,\qquad a_\lambda(G)\ge0.
\]
The supplied registry labels this problem open; that label is supplied curation, not an independently verified literature assessment.
\subsection{Short English statement}
The conjecture asks whether every finite claw-free graph has a chromatic symmetric function whose expansion in the Schur basis has only nonnegative coefficients. It is attributed to Vesselin Gasharov and appears as Stanley's Problem 22. The graph convention and supplied open status require the qualifications stated above.
\subsection{Sources}
\begin{itemize}
\item Local UnsolvedMath registry, permanent ID \(801197\), source code \texttt{catalogue entry}, supplied complete source record. The record lists Stanley as proposer, gives the year 1999, and labels the problem open. Its creation and update timestamps are both \texttt{2026-10-16T14:00:00Z}.
\item Stanley, \emph{Positivity Problems and Conjectures in Algebraic Combinatorics} (1999), Problem 22. Supplied source URL: \url{https://math.mit.edu/~rstan/papers/problems.pdf}. The supplied provenance records access on \texttt{2026-10-16} and extraction method \texttt{pdf+html\_verified}; these are supplied provenance claims, not independent verification performed for this draft.
\item The supplied background attributes the conjecture to Vesselin Gasharov and reports it in Stanley, \emph{Discrete Math.} \textbf{193} (1998), 267--286, Conjecture 1.4. No article title or separate URL for this reference is supplied.
\end{itemize}
\end{document}
